\documentclass{article}
\usepackage{amsmath}
\usepackage{amssymb}
\usepackage{geometry}
\usepackage{enumitem}
\usepackage{array}
\usepackage{titlesec}
\usepackage{xcolor}
\geometry{margin=1in}

% Custom section formatting for hierarchical cases
\titleformat{\subsubsection}[runin]{\normalfont\bfseries}{\thesubsubsection.}{1em}{}
\titleformat{\paragraph}[runin]{\normalfont\bfseries}{\theparagraph.}{1em}{}

% Custom environments for mathematical steps
\newenvironment{stepdetail}
  {\begin{quote}\small\color{blue}}
  {\end{quote}}

\title{Complete 4-Step Variable Elimination Process: \\
No Constraint/Inequality Separation - Pure Inequality System}
\author{Simplified Rational Inequality Solver}
\date{\today}

\begin{document}
\maketitle

\section{Overview}

This document presents a complete 4-step variable elimination process where
\textbf{ALL results are treated as inequalities within cases}. There is
\textbf{no distinction between constraints and inequalities} - everything
must be satisfied within each case.

\textbf{Key Principle:}
\begin{quote}
\textit{At the end, we can just say all the inequalities in one of the cases needs to be true.}
\end{quote}

\textbf{System Features:}
\begin{itemize}
\item No separation between constraints and inequalities
\item All results are inequalities within cases
\item Detailed 4-step process for each variable
\item Complete mathematical documentation
\item Hierarchical case structure
\end{itemize}

\section{Initial System}

Starting with the following system of inequalities involving:
\begin{itemize}
\item \textbf{Variables}: $x_1, x_2, x_3, x_4, \ldots$ (to be eliminated in order)
\item \textbf{Constants}: $a, b, c, d, e, f, g, h, \ldots$ (parameters)
\end{itemize}

\begin{align}
& \text{(1)} \quad x_{1} + x_{2} \leq a \\
& \text{(2)} \quad x_{1} \geq b \\
& \text{(3)} \quad x_{2} \geq c \\
\end{align}

Variables to eliminate (in order): $\{x_{1}, x_{2}\}$

\textbf{Goal}: Eliminate all variables to obtain a system of inequalities
involving only the constants $a, b, c, d, e, f, g, h, \ldots$

\section{Case 1: $x_{1} + x_{2} \leq a$, $x_{1} \geq b$, $x_{2} \geq c$}

Inequalities for Case 1:
\begin{align}
& x_{1} + x_{2} \leq a \\
& x_{1} \geq b \\
& x_{2} \geq c \\
\end{align}

\textbf{Elimination of variable $x_{1}$ in Case 1:}

\textbf{Step 1}: Inequalities after isolating $x_{1}$:
\begin{align}
& x_{1} \leq a - x_{2} \\
& x_{1} \geq b \\
& x_{2} \geq c \\
\end{align}

\textbf{Step 2}: Inequalities after combining to min/max form:
\begin{align}
& x_{2} \geq c \\
& \max\{b\} \leq x_{1} \leq \min\{a - x_{2}\} \\
\end{align}

\textbf{Step 3}: Inequalities after removing $x_{1}$:
\begin{align}
& x_{2} \geq c \\
& b \leq a - x_{2} \\
\end{align}

\textbf{Step 4}: Expanded into 1 subcase(s). Inequalities for each subcase are shown below.

\section{Case 1: $x_{2} \geq c$, $b \leq a - x_{2}$, $- a + b + x_{2} \leq 0$}

Inequalities for Case 1:
\begin{align}
& x_{2} \geq c \\
& b \leq a - x_{2} \\
& - a + b + x_{2} \leq 0 \\
\end{align}

\textbf{Elimination of variable $x_{2}$ in Case 1:}

\textbf{Step 1}: Inequalities after isolating $x_{2}$:
\begin{align}
& x_{2} \geq c \\
& b \leq a - x_{2} \\
& x_{2} \leq a - b \\
\end{align}

\textbf{Step 2}: Inequalities after combining to min/max form:
\begin{align}
& b \leq a - x_{2} \\
& \max\{c\} \leq x_{2} \leq \min\{a - b\} \\
\end{align}

\textbf{Step 3}: Inequalities after removing $x_{2}$:
\begin{align}
& c \leq a - b \\
\end{align}

\textbf{Step 4}: Expanded into 1 subcase(s). Inequalities for each subcase are shown below.

\section{Case 1: $c \leq a - b$, $- a + b + c \leq 0$}

Inequalities for Case 1:
\begin{align}
& c \leq a - b \\
& - a + b + c \leq 0 \\
\end{align}

\section{Case 1: $x_{2} \geq c$, $b \leq a - x_{2}$, $- a + b + x_{2} \leq 0$}

Inequalities for Case 1:
\begin{align}
& x_{2} \geq c \\
& b \leq a - x_{2} \\
& - a + b + x_{2} \leq 0 \\
\end{align}

\textbf{Elimination of variable $x_{2}$ in Case 1:}

\textbf{Step 1}: Inequalities after isolating $x_{2}$:
\begin{align}
& x_{2} \geq c \\
& b \leq a - x_{2} \\
& x_{2} \leq a - b \\
\end{align}

\textbf{Step 2}: Inequalities after combining to min/max form:
\begin{align}
& b \leq a - x_{2} \\
& \max\{c\} \leq x_{2} \leq \min\{a - b\} \\
\end{align}

\textbf{Step 3}: Inequalities after removing $x_{2}$:
\begin{align}
& c \leq a - b \\
\end{align}

\textbf{Step 4}: Expanded into 1 subcase(s). Inequalities for each subcase are shown below.

\section{Case 1: $c \leq a - b$, $- a + b + c \leq 0$}

Inequalities for Case 1:
\begin{align}
& c \leq a - b \\
& - a + b + c \leq 0 \\
\end{align}

\section{Case 1: $c \leq a - b$, $- a + b + c \leq 0$}

Inequalities for Case 1:
\begin{align}
& c \leq a - b \\
& - a + b + c \leq 0 \\
\end{align}

\section{Final Solution Summary}

After systematically eliminating all variables $\{x_{1}, x_{2}\}$
using the enhanced hierarchical 4-step process, we obtain the following
1 final case(s), each representing a complete solution region:

\subsection{Final Case 1: Case 1}

This case represents the solution region defined by:
\begin{align}
& c \leq a - b \\
& - a + b + c \leq 0 \\
\end{align}

\begin{stepdetail}
\textbf{Interpretation}: This final case represents all parameter values
$(a, b, c, d, e, f, g, h, \ldots)$ for which the original system has solutions.
\end{stepdetail}

\section{Conclusion}

The complete 4-step variable elimination process with no constraint/inequality separation provides:
\begin{itemize}
\item \textbf{Unified Treatment}: All results are inequalities within cases
\item \textbf{Mathematical Clarity}: No artificial distinction between constraint types
\item \textbf{Systematic Elimination}: Variables are eliminated in specified order
\item \textbf{Hierarchical Organization}: Complex case structures are clearly organized
\item \textbf{Simple Interpretation}: Final result requires all inequalities in one case to be true
\end{itemize}

This elimination process demonstrates that:
\begin{itemize}
\item All results are expressed as inequalities within cases
\item No distinction is made between constraints and inequalities
\item The final solution requires all inequalities in one of the cases to be true
\item The 4-step process systematically eliminates each variable
\end{itemize}

The final result characterizes the complete solution space through
1 distinct case(s), where each case contains only inequalities
that must all be satisfied simultaneously.

\end{document}